\documentclass[a4paper,10pt]{article}

\usepackage[utf8]{inputenc}
\usepackage[absolute]{textpos}
\usepackage{changepage}
\usepackage{color,soul}
\usepackage{ulem}
\usepackage{blindtext}
\usepackage{rotating}
\usepackage{nopageno}
\usepackage{circledsteps}
\usepackage{latexsym}
\usepackage{parskip}
\usepackage{stix}
\usepackage{tikz}
\usepackage{xcolor}
\usepackage[paperwidth=5.5in, paperheight=20in, margin=0.7in]{geometry}
\usetikzlibrary{arrows.meta}
%dimensions 5in by 7.75; I have altered to bring in line with p.20 layout]
\renewcommand{\labelitemi}{-}
\setulcolor{red}

\begin{document}
\color{blue}
\begin{flushright}
\Circled{Notes \& orig. drafts (?) page 1}\par
\end{flushright}
\begin{flushleft}
\color{red} 
\begin{turn}{25}%
* See \ul{Taylor p215} 
\end{turn}
\ul{80 Flowers} (possible substitutes\par
\color{red}
for the 80 listed on ppI-VI) \hfill \color{blue} \Circled{{$\stackrel{\hbox{orig p \#}}{\hbox{IA-B}}$}}\par
\color{red}
\small
% line (join "See" and 3)
\begin{tikzpicture}[remember picture,overlay]
    \color{red}
    \draw (0, 0) .. controls (-1, 1) .. (0, 2);
\end{tikzpicture}
\Circled{3} perennial .. weedy .. Eurasian .. poppy family .. {$\stackrel{\hbox{called also}}{\hbox{killwort sightwort}}$}\par
%above lines not great at the moment in making clear where LZ starts new line for space, squeezes in a bit on the right etc
\Circled{2} the lesser celandine or pilewort, still used in home med. for piles\par
``occurs locally in Eastern states.\par
\begin{itemize}
\color{blue}
\setulcolor{blue}
\normalsize
\item Celandine, \ul{Chelidonium majus} (\ul{Kamm}\par
\tiny 
\color{red}
swallowort - Lewis \& Short\par
\color{blue}
\normalsize
\color{red}
\Circled{$\stackrel{\hbox{{*See:}}}{\hbox{\color{blue}{pp.40-42}}}$} \color{blue}Port Jefferson planted Oct 18/74)\par
\color{red}
\small
\Circled{1} (Wordsworth's
\begin{tikzpicture}[remember picture,overlay]
    \color{red}
    \draw[->] plot [smooth] coordinates {(-1.5, -0.1) (-4, 0) (-4, 2) (-2.8, 3)};
\end{tikzpicture}
 poems refer to \ul{Ranunculus ficaria}, a crowfoot.\par
% tick
\begin{tikzpicture}[remember picture,overlay]
    \color{red}
    \draw (1.2, 1) --  (2, -0.5) -- (4.4, 7);
\end{tikzpicture}
see \# 65 vi pV 
%entry for 65 is on pV of working notebook, but gives sage
\color{blue}
\hfill
$\uparrow$ L = little frog $\stackrel{\hbox{Buttercup}}{\hbox{i.e. meadow habit}}$
%arrow should point to Ranunculus
\rule{10cm}{0.01cm}
\normalsize
\item Lavender cotton, \ul{Santolina chamaecypar-} \par
\ul{-issus} \ul{Kamm} p 206 illust \#25 Pt. J, planted\par 
%looks v like a new entry for -issus but clearly cont from previous line
\color{red}
\tiny
\newsavebox{\abc}
\savebox{\abc}{
    \fcolorbox{blue}{white}{
        \begin{minipage}{0.5in}
            How great fallen are the great L.Z.
        \end{minipage}
    }
}
\begin{turn}{25}\usebox{\abc}\end{turn}
%a 2 and 3 in there? presumably an LZ quotation
\color{blue}
\normalsize
Oct 18/74
\small
\color{red}
John Parkinson 1629 \ul{Paradisus Terrestis} ``The rarity \& novelty of this herb, being for the most part but in the gardens of great persons, doth cause it to be of great regard"
%LZ tweaking the quotation?
\color{blue}
\normalsize
\item Spearmint
\begin{tikzpicture}[remember picture,overlay]
    \color{red}
    \draw (-0.8,0.3) --  (-0.5, -0.1) -- (0.5, 0.5);
\end{tikzpicture}
(Pt. J. 
\begin{tikzpicture}[remember picture,overlay]
    \color{red}
    \draw (-0.8,0.3) --  (-0.5, -0.1) -- (0.5, 0.5);
\end{tikzpicture}
planted 
\begin{tikzpicture}[remember picture,overlay]
    \color{red}
    \draw (-0.8,0.3) --  (-0.5, -0.1) -- (0.5, 0.5);
\end{tikzpicture}
in Oct 74)
\begin{tikzpicture}[remember picture,overlay]
    \color{red}
    \draw (-0.8,0.3) --  (-0.5, -0.1) -- (0.5, 0.5);
\end{tikzpicture}
\item Sp\ul{ider}
\begin{tikzpicture}[remember picture,overlay]
    \color{red}
    \draw (-0.8,0.3) --  (-0.5, -0.1) -- (0.5, 0.5);
\end{tikzpicture}
pl\ul{ant} 
\begin{tikzpicture}[remember picture,overlay]
    \color{red}
    \draw (-0.8,0.3) --  (-0.5, -0.1) -- (0.5, 0.5);
\end{tikzpicture}
(lily family)
\hspace{0.5in}
{$\stackrel{\hbox{\tiny{i.e. LZ How great are the fallen}}}{\hbox{}}$}
\begin{tikzpicture}[remember picture,overlay]
    \color{blue}
    \draw[->] (-5.5, 0.3) .. controls (-3, 0.5) .. (-2.5, 0.2);
\end{tikzpicture}
\item \ul{thyme (time)} (Pt. J. {$\stackrel{\hbox{Aug-}}{\hbox{Sep}}$}74) (see \#70 pVI v.i)
\item (\ul{mint} 
\begin{tikzpicture}[remember picture,overlay]
    \color{blue}
    \draw (-0.8,0.3) --  (-0.1, -0.6);
    \draw (-0.1, 0.3) -- (-0.8,-0.6);
\end{tikzpicture}
\ul{geranium} 
\begin{tikzpicture}[remember picture,overlay]
    \color{blue}
    \draw (-0.8,0.3) --  (-0.1, -0.6);
    \draw (-0.1, 0.3) -- (-0.8,-0.6);
\end{tikzpicture}
or costmary 
\begin{tikzpicture}[remember picture,overlay]
    \color{blue}
    \draw (-0.8,0.3) --  (-0.1, -0.6);
    \draw (-0.1, 0.3) -- (-0.8,-0.6);
\end{tikzpicture}Pt. J 
\begin{tikzpicture}[remember picture,overlay]
    \color{blue}
    \draw (-0.8,0.3) --  (-0.1, -0.6);
    \draw (-0.1, 0.3) -- (-0.8,-0.6);
\end{tikzpicture}
Oct 
\begin{tikzpicture}[remember picture,overlay]
    \color{blue}
    \draw (-0.8,0.3) --  (-0.1, -0.6);
    \draw (-0.1, 0.3) -- (-0.8,-0.6);
\end{tikzpicture}
26/74
\begin{tikzpicture}[remember picture,overlay]
    \color{blue}
    \draw (-0.8,0.3) --  (-0.1, -0.6);
    \draw (-0.1, 0.3) -- (-0.8,-0.6);
\end{tikzpicture}\par 
Chrysanthemum Balsomita Garden Ency Taylor)
\begin{tikzpicture}[remember picture,overlay]
    \color{blue}
    \draw (-6, 0.3) --  (-5.5, -0.5) -- (0, 14);
\end{tikzpicture}
\begin{adjustwidth}{-0.5in}{-0.5in}
    \rule{4.5in}{0.005in}  
\end{adjustwidth}
%this item crossed out; probably Balsamita, but doesn't quite look like it; line down to
\begin{flushright}
% may want an alternative arrow to make clear connection / complete line to material above    
$\downarrow$(identified as follow at {$\stackrel{\hbox{Suffolk}}{\hbox{Rose Soc.}}$})
\end{flushright}
\item sweet-fern (Comptonia, \ul{Taylor}; resembles spleenwort (a fern) but \ul{sweet-fern} is a shrub of the walnut family - see Gray)
Oct 28/74 Pt..J.\par
\color{red}
\footnotesize
\begin{minipage}{2in} 
    Taylor lists wild geranium under\par
    \ul{geraniaceae} -/a separate\par
    family
\end{minipage}%
\hfill
\begin{minipage}{0.6in}
    Gray $\stackrel{\hbox{817-}}{\hbox{p818}}$ 
\end{minipage}%
\begin{minipage}{0.7in}
    under\par
    g\'eum\par
    (Rose family)
\end{minipage}%
\begin{tikzpicture}[remember picture,overlay]
    \color{red}
    \draw[->] (-1.8,-0.3) -- (-4, -1);
\end{tikzpicture}
%Gray has Geum with grave, but LZ's looks acute
\color{blue}
\normalsize
\item wild geranium (\ul{Taylor}, \ul{Ge}\setulcolor{red}r\ul{anicum} \ul{macu}la-\par
\ul{tum}; also called alumsroot
\color{red}
* 
\color{blue}
\& chocolateflower\par
\setulcolor{blue}
\color{red}
\begin{textblock*}{3.22in}(0.5cm,20.5cm)%
    \small
    \begin{minipage}{3.22in} 
        \color{red}
        \setulcolor{red}
        but $\stackrel{\hbox{is}}{\hbox{$\caretinsert$}}$ not\par
        \ul{Heuchera}\par
        also called\par
         alumroot\par
        (saxifrage\par
        family)
        \normalsize
  \end{minipage}%
  \end{textblock*}%
\color{blue}
\setstcolor{red}
[\st{the relation to those}]
\begin{tikzpicture}[remember picture,overlay]
    \color{red}
    \draw[->] (-0.1, 0.1) .. controls (-1.5, 0) .. (-3, 0.25);
\end{tikzpicture}
{$\stackrel{\hbox{identified*}}{\hbox{Oct 30/74}}$}\par
*neighbor called it wild strawberry, but\par 
no berry\par
%should have square brackets from [*neighbor to berry] but was confusing LaTeX
Gray - The name of a nymph etc orig a water plant\par 
sour gum, pepperidge \color{red} (\ul{Cheyney} What Tree is That?)
\begin{textblock*}{2in}(0.8cm,23.5cm)%
    \small
    \begin{minipage}{2in} 
        \color{red}
        See\par
        Burpee\par
        \sout{???}\par
        \normalsize
  \end{minipage}%
  \end{textblock*}%
\color{blue}
\item Nyssa sylvatica Pt J 306 {$\stackrel{\hbox{- identified?}}{\hbox{E Brway Nov 3/74?}}$}
\begin{adjustwidth}{-0.5in}{-0.5in}
    \rule{4.5in}{0.005in}  
\end{adjustwidth}


\item Catnip - Same as \underline{Glecoma} hederacea (Taylor)\par
\setulcolor{red}
?\#29 p9 III v.i. ``The \ul{catnip and the} \par
\ul{amaranth! - man's} \ul{earthly household} \par
\ul{peace} \& \ul{the ever-encroa}ch\ul{ing appetite} \ul{for God}"\par
\setulcolor{blue}
Melville, \ul{Pierre} Bk xxv IV (p356 \ul{Signet} ed.)
\color{red}
\begin{textblock*}{2in}(0.5cm,26cm)%
    \small
    \begin{minipage}{2in} 
        \color{red}
        unwithering\par
        unfading $\rightarrow$
        \normalsize
  \end{minipage}%
  \end{textblock*}%

\color{red}
\begin{adjustwidth}{-0.5in}{-0.5in}
catnip - N\'epeta %more accents 
Catmint .. per. \& ann. herbs .. \sout{tall}
\sout{to erect or trailing} tall \& erect, or dwarf trailing, generally (over)
\end{adjustwidth}
\end{itemize}
\end{flushleft}


\end{document}